\scanpage{257}
\mspage{207}

qu'on peut définir via la comparaison de la ``théorie $C$'' d'une part, la ``théorie $C'$'' de l'autre, avec ce que nous donne directement la flèche composée \unsure{(x)|(*)}. (p. 204).

Finalement, il me semble que l'\uline{existence} d'une théorie des ``fibres homotopiques'' (\unsure{i.e.|. . .} ``$W$-homotopiques'') d'une flèche $X\to S$ dans $\mathrm{Cat}$, résulte de l'équivalence \unsure{(xx)|(**)} (p. 204), et \inserted{cette théorie} pourrait sans doute dans une certaine mesure se développer à partir de là, sans choix d'un foncteur $C$ ni d'une structure de ``catégorie de modèles au sens de Quillen\unsure{''|,}

\unsure{9.|4.} \uline{Axiomatisation des ``fibres \inserted{$W$-}homotopiques.''} La réflexion de hier me fait revenir sur l'axiomatique dans des ``cat. de modèles''
\[
\unsure{(\,.\quad {}'|(\,.\quad .}\qquad(M,W)
\]
% slanted
\margin{$W$ satisfait l'\unsure{``|{}}axiome de saturation faible~: les iso dans $W$, et si deux parmi $f,g,gf$ dans $W$, le troisième \unsure{itou|aussi} \unsure{(si)|(sic)|bien sûr)}.}
plus générales que $\mathrm{Cat}$, du $1^{\mathrm{er}}$ \unsure{append.} de XII (\unsure{cf th.-scholie 3|cf th.-scholie 2} p. 108--113) Pour le moment, on va supposer seulement $M$ stable par produits fibrés (on n'a pas besoin d'objet final). On introduit l'ensemble \deleted{\unsure{des} $\mathrm{Fib}_W\subset\mathrm{Fl}(M)$}
\[
\mathrm{Fib}_W\subset\mathrm{Fl}(M),
\]
% archive sheet number: 257
